%% parts/p05/decidable-arithmetics.tex --- Part 5
\chapter{Decidable Arithmetics}\label{sec:p05-decidable-arithmetics}

\noindent\textit{Reader note.}  This separately authored comparison follows
the exact inherited undecidability results.  It records decision procedures
for two proper reducts and does not weaken or replace the inherited theorem for
full arithmetic.  The historical and method sources are
\cite{Presburger1929,Skolem1930Arithmetic,FefermanVaught1959}.

The undecidability of arithmetic with both addition and multiplication does not
transfer to each reduct. This chapter fixes two classical reducts and states
exactly what their decision procedures preserve.

\begin{hbdefinition}{def:p05-presburger}{Presburger arithmetic}
Presburger arithmetic is the complete first-order theory of
\((\mathbb N;0,1,+,<)\). Congruence predicates
\(x\equiv r\pmod m\) may be added as definitional expansions for quantifier
elimination. The name is necessary to distinguish this additive theory from
Peano arithmetic and from quantifier-free linear integer arithmetic. It stands
as a complete theory of the standard additive structure, differs by omitting
multiplication of variables, and is used in the elimination theorem
\cite{Presburger1929}.
\index{Presburger arithmetic}
\end{hbdefinition}

\begin{hbtheorem}{thm:p05-presburger-qe}{cited=Presburger1929}{Quantifier elimination and decidability}
After adding all congruence predicates, every Presburger formula is equivalent
over \(\mathbb N\) to a quantifier-free formula. Consequently Presburger
arithmetic is complete and decidable.
\end{hbtheorem}

An elimination step first normalizes coefficients and inequalities, replaces a
quantified variable by finitely many residue cases modulo the least common
multiple of relevant coefficients, and checks boundary cases determined by
the remaining linear terms. The finite residue split establishes termination;
it can cause severe formula growth. Fischer and Rabin prove a doubly
exponential lower bound in formula length in their computation model
\cite{FischerRabin1974}. Decidability therefore does not imply practical
tractability.

\begin{hbexample}{def:p05-presburger-example}{A congruence elimination}
The formula \(\exists y\,(x=3y+1)\) is equivalent over \(\mathbb N\) to
\(x\equiv1\pmod3\). The congruence predicate records the residue information
that would be lost by trying to eliminate \(y\) in the language of order and
addition without a definitional expansion.
\end{hbexample}

\begin{hbdefinition}{def:p05-skolem-arithmetic}{Skolem arithmetic}
Skolem arithmetic is the complete first-order theory of
\((\mathbb N_{>0};1,\times)\). The positive-domain convention avoids a special
absorbing-zero case; including zero gives an interdefinable variant only after
that case is handled explicitly. The name is needed to distinguish the
multiplicative reduct from full arithmetic. It stands as a complete theory,
differs by omitting addition and order, and is used in the prime-exponent
decision argument \cite{Skolem1930Arithmetic}.
\index{Skolem arithmetic}
\end{hbdefinition}

Unique factorization sends a positive integer to its finitely supported vector
of prime exponents. Multiplication becomes coordinatewise addition, while
quantification ranges over finite-support vectors. Feferman--Vaught-style
decomposition separates the uniform coordinate theory from finite support and
yields decidability. This is not a definitional equivalence with one fixed
finite power of Presburger arithmetic: the set of prime coordinates is
unbounded \cite{FefermanVaught1959}.

\begin{hbtheorem}{thm:p05-skolem-decidable}{cited=FerranteRackoff1979}{Decidability of Skolem arithmetic}
The first-order theory of \((\mathbb N_{>0};1,\times)\) is decidable.  Rackoff's
weak-direct-power procedure decides sentences of length \(n\) in space
\(2^{2^{2^{cn}}}\) for some constant \(c\).  Thus the published upper bound is
elementary recursive, although it is neither polynomial nor a claim of
practical tractability \cite{Rackoff1976}.
\end{hbtheorem}

The contrast is exact: the additive and multiplicative reducts are decidable,
whereas the combined standard structure \((\mathbb N;0,1,+,\times)\) interprets
the coding used by the incompleteness and undecidability theorems. A decision
result for a reduct cannot be transferred after restoring the omitted
operation.

\begingroup\small
\begin{longtable}{@{}p{0.20\textwidth}p{0.25\textwidth}p{0.25\textwidth}p{0.20\textwidth}@{}}
\caption{Decidable arithmetic theories.  The complexity
entries are published worst-case bounds, not matching upper and lower bounds.}
\label{tab:p05-decidable-arithmetics}\\
\toprule
Theory and language & Decision method & Published complexity boundary &
What makes the reduct proper \\
\midrule
\endfirsthead
\toprule
Theory and language & Decision method & Published complexity boundary &
What makes the reduct proper \\
\midrule
\endhead
\(\operatorname{Th}(\mathbb N;0,1,+,<)\) &
Quantifier elimination after adding all congruence predicates; finite residue
splitting. &
Fischer--Rabin: a doubly exponential lower bound in formula length in their
computation model. &
No multiplication of two variables. \\
\addlinespace
\(\operatorname{Th}(\mathbb N_{>0};1,\times)\) &
Prime-exponent vectors and weak-direct-power decomposition with finite support. &
Rackoff: space \(2^{2^{2^{cn}}}\) for some constant \(c\) in the stated
procedure. &
No addition, order, or absorbing zero. \\
\addlinespace
\(\operatorname{Th}(\mathbb N;0,1,+,\times)\) &
No decision procedure: theoremhood is undecidable. &
Not a complexity upper bound; the decision problem is undecidable. &
Both operations are present and support the fixed arithmetization. \\
\bottomrule
\end{longtable}
\endgroup
